
% =====================================================================
\section{Three derivations of the floor}
\label{sec:floor}
% =====================================================================

The Boundary--Thickness Theorem produced the floor $\thick>0$ from the geometry
of separators. We now give two further, independent derivations---one
representational, one cardinality-theoretic---and a cost argument showing that
driving a separator toward zero content has diverging cost. The three agree:
each bounds the same constant from below away from zero, and none can be reduced
to the others. This redundancy is the paper's evidence that the floor is
structural and not an artefact of one formalism.

\subsection{The representational derivation}

The first derivation asks what it costs to \emph{name} a region, independently of
any separator. An embedded agent identifies a region by a finite descriptor
drawn from a finite scheme; we show no such scheme distinguishes all regions of
a bounded resolvable space exactly, so a positive descriptive residue persists.

\begin{definition}[Descriptive scheme]
\label{def:scheme}
A \emph{descriptive scheme} for a partition $\Part$ with $N$ cells is an
injection $\rho$ from a set $\mathcal{R}$ of realisable regions into the finite
set $\{0,1\}^{N}$ of cell-indicator strings: $\rho(A)$ is the string whose
$i$-th entry is $1$ iff $\cell_i\subseteq A$. A region $A$ is
\emph{exactly named} by $\rho$ if $A=\bigcup\{\cell_i:\rho(A)_i=1\}$.
\end{definition}

\begin{proposition}[Representational residue]
\label{prop:rep-residue}
Under \Cref{ax:brs}, for any region $A$ that is not $\Part$-realisable the
descriptive scheme $\rho$ names instead the realisable region
$\widehat A=\bigcup\{\cell_i:\cell_i\subseteq A\}$, and the discrepancy
$A\,\triangle\,\widehat A$ lies in the separator $\Sigma_\Part(A)$, of measure
$\ge\thick$. No finite descriptive scheme over $\Part$ names every region of
$\Space$ exactly; the unnamed discrepancy has measure at least the floor.
\end{proposition}

\begin{proof}
A scheme over $\Part$ has range $\{0,1\}^N$, a finite set, while the regions of
$\Space$ (Borel sets modulo null sets) are uncountable under \ref{brs:hier}; so
$\rho$ cannot be injective on all regions and must collapse some region $A$ onto
a realisable $\widehat A$. The cells contributing to $A$ only partially lie in
$\Sigma_\Part(A)$: a cell $\cell_i$ with $\cell_i\subseteq A$ is named, a cell
disjoint from $A$ is omitted, and a cell meeting both $A$ and $\Space\setminus A$
is exactly a separator cell. Hence $A\,\triangle\,\widehat A\subseteq
\Sigma_\Part(A)$, and where $A$ is non-realisable this symmetric difference is
non-null and contained in a separator of measure $\ge\thick$ by
\Cref{thm:thickness}.
\end{proof}

\begin{remark}[A faithful image is not the invariant]
\label{rem:faithful}
\Cref{prop:rep-residue} says that the best name a finite scheme can give a region
is a realisable approximant differing from it by at least a floor's worth of
content. A representation that is faithful at the resolution of $\Part$ is still
not the region itself; the gap is the same $\thick$ that the separator carries.
The descriptive and the geometric residues coincide.
\end{remark}

\subsection{The cardinality derivation}

The second derivation needs no geometry at all---only counting. It shows that the
collection of distinctions an agent might draw outruns any finite stage of
representation, so that at each finite stage a positive shortfall remains.

\begin{proposition}[Cardinality shortfall]
\label{prop:card}
Under \ref{brs:hier}, the family of regions of $\Space$ \textup{(}Borel modulo
null\textup{)} has cardinality of the continuum, while the regions exactly
nameable at any finite stage $\Part_n$ number at most $2^{N(\Part_n)}<\infty$.
Hence at every finite stage the nameable regions are a measure-dense but
cardinality-negligible subfamily, and the act of naming any further region
requires passing to a strictly finer partition, which by \Cref{prop:refine}
carries its own positive floor $\thick_{n+1}>0$.
\end{proposition}

\begin{proof}
By \Cref{lem:finite-cells} each $\Part_n$ is finite, so its nameable regions
(unions of cells) number $2^{N(\Part_n)}$. By \ref{brs:hier} the partitions
generate the Borel $\sigma$-algebra modulo null sets, whose cardinality is that
of the continuum; no finite stage attains it. To name a region not realisable at
stage $n$ one refines to a stage $m>n$ at which it becomes a union of cells; by
\Cref{prop:refine} every such stage carries $\thick_m\ge\mu_{\min}(\Part_m)>0$.
Thus the floor reappears at each refinement; there is no finite stage of exact,
complete naming and no stage with zero floor.
\end{proof}

\begin{remark}
\Cref{prop:card} is a pure counting fact and is logically independent of the
measure-geometry of \Cref{sec:thickness}: even if every separator could somehow
be made thin, the finite nameable family could never exhaust the continuum of
regions, so a fresh refinement---with its own floor---would always be required.
The two derivations corner the floor from disjoint directions.
\end{remark}

\subsection{The cost derivation}

The third derivation treats the floor as the limit of an optimisation: it shows
that the cost of thinning a separator toward zero content diverges, so a finite
agent never reaches the sharp cut. We work with an abstract, monotone cost
functional and assume only what any reasonable accounting of refinement must
satisfy.

\begin{definition}[Refinement cost]
\label{def:cost}
A \emph{refinement cost} is a functional $\Gamma$ assigning to each realisable
partition $\Part$ a value $\Gamma(\Part)\in[0,\infty)$ such that
\begin{enumerate}[label=\textup{(C\arabic*)},leftmargin=3em]
\item\label{c:mono} \textbf{Monotone.} $\Part'\preceq\Part\Rightarrow
\Gamma(\Part')\ge\Gamma(\Part)$ \textup{(}finer costs at least as much\textup{)};
\item\label{c:additive} \textbf{Additive over cells.} splitting a cell of
measure $m$ into cells of measures $m_1,\dots,m_r$ adds
$\sum_s g(m_s)-g(m)\ge0$ to $\Gamma$, for a convex $g\colon(0,\infty)\to\R$
with $g(m)\to+\infty$ as $m\to0^+$;
\item\label{c:reso} \textbf{Resolution-bounded.} achieving boundary thickness
$\thick_\Part(A)=t$ for a fixed region requires a partition whose separating
cells have measure $\le t$, hence each of cost $\ge g(t)$.
\end{enumerate}
\end{definition}

\begin{theorem}[Cost divergence at the sharp cut]
\label{thm:cost}
Let $\Gamma$ be a refinement cost. For a fixed region $A$, the cost of realising
boundary thickness $t$ satisfies $\Gamma_t\ge g(t)\to+\infty$ as $t\to0^+$.
Hence no finite-cost separation process attains a sharp ($t=0$) separator; the
infimum of attainable thickness over finite-cost processes is a strictly
positive $\thick$.
\end{theorem}

\begin{proof}
By \ref{c:reso}, realising thickness $t$ requires separating cells of measure
$\le t$; by \ref{c:additive} each such cell costs at least $g(t)$, and since at
least one separating cell is present (\Cref{lem:nonempty-sep}),
$\Gamma_t\ge g(t)$. Since $g$ is convex with $g(m)\to+\infty$ as $m\to0^+$,
$\Gamma_t\to+\infty$ as $t\to0^+$. A finite-cost process has $\Gamma\le C<\infty$
for some bound $C$, so its attainable thickness is bounded below by
$t_\ast:=\sup\{t:g(t)\le C\}>0$. Setting $\thick=t_\ast$ gives a strictly
positive floor below which no finite-cost process reaches.
\end{proof}

\begin{remark}[Bounded return forces a halt]
\label{rem:bounded-return}
The cost derivation has a corollary used in \Cref{sec:inquiry}: because the
return from thinning a separator---the additional content correctly assigned---is
bounded by the finite measure $\mu(\Space)$ (\ref{brs:bdd}), while the cost
$g(t)$ diverges as $t\to0^+$, the marginal return per unit cost falls to zero.
A finite agent with any positive cost-per-return threshold therefore halts at a
strictly positive thickness. The halt is not a failure of resolve but the
rational endpoint of a diverging-cost, bounded-return optimisation.
\end{remark}

\subsection{Agreement of the three}

\begin{theorem}[The floor is one constant]
\label{thm:floor-agree}
For a fixed bounded resolvable space and agent, the geometric floor of
\Cref{def:floor}, the representational residue of \Cref{prop:rep-residue}, and
the cost-limited thickness of \Cref{thm:cost} are all strictly positive and
bound the same quantity: each equals or exceeds $\mu_{\min}$ at the agent's
finest realisable partition, and each vanishes only in the excluded limit
$\reso\to0$. No one of them can be reduced to zero while the others remain
positive.
\end{theorem}

\begin{proof}
All three are controlled by $\mu_{\min}$ at the finest realisable partition. The
geometric floor satisfies $\thick\ge\mu_{\min}$ (\Cref{thm:thickness}); the
representational residue is a separator of measure $\ge\thick\ge\mu_{\min}$
(\Cref{prop:rep-residue}); the cost-limited thickness $t_\ast$ satisfies
$g(t_\ast)\le C$ with separating cells of measure $\ge\mu_{\min}$ forced by
\ref{brs:res}, so $t_\ast\ge\mu_{\min}$. Each is therefore $\ge\mu_{\min}>0$, and
each $\to0$ only as $\mu_{\min}\to0$, i.e.\ $\reso\to0$, which \ref{brs:res}
excludes for an embedded agent. Hence they rise and fall together and are
positive together.
\end{proof}
